% Generated draft: requires mathematical and source review.
\documentclass[11pt]{article}
\usepackage[margin=25mm]{geometry}
\usepackage{fontspec}
\setmainfont{Latin Modern Roman}
\usepackage{amsmath,amssymb,mathtools,unicode-math}
\setmathfont{Latin Modern Math}
\usepackage{enumitem,needspace,xurl,hyperref}
\setlength{\parindent}{0pt}
\setlength{\parskip}{5pt}
\setlength{\emergencystretch}{3em}
\newcommand{\N}{\mathbb{N}}
\newcommand{\Z}{\mathbb{Z}}
\newcommand{\Q}{\mathbb{Q}}
\newcommand{\R}{\mathbb{R}}
\newcommand{\C}{\mathbb{C}}
\newcommand{\F}{\mathbb{F}}
\newcommand{\A}{\mathbb{A}}
\newcommand{\PP}{\mathbb{P}}
\newcommand{\E}{\mathbb{E}}
\newcommand{\eps}{\varepsilon}
\newcommand{\dd}{\,\mathrm{d}}
\begin{document}
\section*{UnsolvedMath ID: 11000010}
\textbf{Source problem code:} catalogue entry
\subsection{Definitions and mathematical statement}
Fix the genus \(g\) and the preceding source's representation space \(X_g(\Gamma)\) and volume maps \(\operatorname{Vol}_k\). For each integer \(1\leq k\leq3g-3\), determine
\[
 \mathcal V_{g,k}(\Gamma)
   =\{\operatorname{Vol}_k(\rho):\rho\in X_g(\Gamma)\}\subset\R,
 \qquad
 \bigcup_{\Gamma\ \mathrm{finitely\ presented}}
         \mathcal V_{g,k}(\Gamma).
\]
The first request fixes \(\Gamma\); the second varies over all finitely presented groups. Also determine the images when one restricts to representations having holomorphic representatives or representatives totally geodesic with respect to a fixed metric, and how those geometric conditions constrain possible volumes.

The excerpt omits the definitions of \(X_g(\Gamma)\), \(\operatorname{Vol}_k\), the geometric representatives and their domains. Consequently no integral, normalization or homology pairing is substituted for the named volume map. The source notes Mirzakhani's recursive Weil--Petersson volume formulas for moduli of surfaces with nonempty totally geodesic boundary; these are contextual moduli-space volumes, not the requested representation-volume image.

A following discussion proposes further computations from linear representations. In the usual conjugacy-class interpretation of \(X_g(\Gamma)\), for \(\rho:\Gamma\to\operatorname{Mod}_g\), \(\psi:\operatorname{Mod}_g\to\operatorname{GL}_m(\C)\), a conjugation-invariant function \(c\) and a fixed \(\gamma\in\Gamma\), composition gives the invariant
\( [\rho]\mapsto c(\psi(\rho(\gamma)))\).
The source also proposes subgroup commensurability invariants obtained from such representations, cites computations in [Su2], and calls for more. It reports that no faithful linear representation was known for \(g\geq3\); this is retained as a dated source assertion, not a current literature conclusion. The construction of commensurability invariants is not supplied.
\subsection{Short English statement}
Determine the images of each representation-volume map, their union over finitely presented groups, and the restrictions imposed by holomorphic or totally geodesic representatives; the volume definitions are missing from the extract.
\subsection{Sources}
\begin{itemize}
\item \emph{Problems on Mapping Class Groups and Related Topics}, edited by Farb, Problem 2.3 (``Volume spectrum''), supplied PDF page 23, with the following linear-representation discussion. \url{https://www.math.uchicago.edu/~farb/papers/mcgbook.pdf}. Recorded access 2026-07-29; extraction pdf-font-structure. The supplied running heading names ``Some problems on mapping class groups and moduli space.' '
\item The source cites Mirzakhani [Mir] for boundary Weil--Petersson recursion and [Su2] for linear-representation computations; complete bibliography entries are absent.
\item Local UnsolvedMath registry, \texttt{data/unsolvedmath/problems.json}, permanent ID \(11000010\), exact source code \texttt{catalogue entry}; supplied OPEN-TRIAGE summary distinguishes the representation-volume question from Mirzakhani's recursion and reports no resolution.
\end{itemize}
\end{document}
